\chapter{Lexicon}
\rule{\textwidth}{0.5pt}    % Line 
\thispagestyle{fancy}


This lexicon provides key terms used throughout the text, presented in a format designed to clearly distinguish the original language, the English translation, and the definition. Each entry consists of six elements:
\begin{itemize}
	\item \textbf{Language Symbol}: A single-letter symbol indicating the language of the term. 
	\begin{itemize}
		\item \textbf{G} for Greek.
		\item \textbf{H} for Hebrew.
		\item \textbf{A} for Aramaic.
		\item \textbf{E} for English.
	\end{itemize}
	\item \textbf{Strong's Reference}: The reference number from Strong's Concordance.
	\item \textbf{Original Word}: The term in its source language. Typically an English transliteration.
	\item \textbf{Pronunciation}: How to pronounce the word using English phonetics.
	\item \textbf{English Translation(s)}: The equivalent word or phrase in English, shown in italics. This can also be a combination of various words this word is translated to as well as how often it is translated to such words.
	\item \textbf{Definition}: An explanation or context of the term.
	\item \textbf{References}: The references used for determining the information.
\end{itemize}

\medskip
To improve readability, each entry is fully color-coded according to the language of the original word. The colors as well as the entry explanations are listed below:
\begin{itemize}
    \item \lexgreek{Strong's number}{word}{pronunciation}{english}{definition}{[references]}
	\item \lexhebrew{Strong's number}{word}{pronunciation}{english}{definition}{[references]}
	\item \lexaramaic{Strong's number}{word}{pronunciation}{english}{definition}{[references]}
	\item \lexenglish{word}{definition}{[references]}
\end{itemize}


\noindent\rule{\textwidth}{0.5pt}
\vspace{0.5em} % vertical space


\lexhebrew{H2616}{ḥesed}{kheh'-sed}{KJV translations: mercy (149x), kindness (40x), loving-kindness (30x), goodness (12x), kindly (5x), merciful (4x), favor (3x), good (1x), goodliness (1x), pity (1x), reproach (1x), wicked thing (1x)}{Kindness; by implication (towards God) piety; rarely (by opposition) reproof, or (subjectively) beauty:—favour, good deed(-liness, -ness), kindly, (loving-) kindness, merciful (kindness), mercy, pity, reproach, wicked thing.}{\cite{BlueLetterLexicon}}

